\chapter{Express}

Abbiamo sviluppato la stessa applicazione tre volte: con CGI e Bash, con PHP e
con Node.js puro. Ogni volta abbiamo incontrato gli stessi punti dolenti:

\begin{itemize}
  \item molto lavoro ripetitivo e noioso (analizzare i cookie, leggere e
        interpretare il corpo delle richieste);
  \item nessuna gestione delle sessioni pronta all'uso (tranne che in PHP);
  \item nessun supporto integrato ai \textit{middleware}, per esempio per
        l'autenticazione;
  \item nessun motore di template integrato.
\end{itemize}

Gestire il routing, generare template, gestire sessioni, autenticazione e
autorizzazione sono compiti \textbf{comuni a ogni progetto}. Non vogliamo
reinventare la ruota ogni volta.

\section{Che cos'è un framework}

\dfn{Framework}{
  Un \term{framework} è un insieme predefinito di componenti software, strumenti
  e buone pratiche che fa da fondamento per lo sviluppo di software. I
  \term{web framework} sono progettati specificamente per semplificare e
  snellire lo sviluppo di applicazioni web, fornendo un modo
  \textbf{strutturato} di costruirle e organizzarle.}

Framework e \textbf{librerie} non sono la stessa cosa, e la differenza è
sostanziale.

\begin{center}
\renewcommand{\arraystretch}{1.4}
\begin{tabular}{@{}>{\raggedright\arraybackslash}p{6.3cm}@{\hspace{1.0cm}}>{\raggedright\arraybackslash}p{6.3cm}@{}}
  \textbf{Libreria} & \textbf{Framework} \\
  È essenzialmente un insieme di funzioni che possiamo invocare per far
  svolgere un compito.
  & Incarna un progetto astratto, con molto più comportamento già
    incorporato. \\
  Ogni chiamata svolge il proprio lavoro e poi \textbf{restituisce il controllo}
  a chi ha chiamato.
  & Il codice scritto da noi viene \textbf{innestato} dentro il framework. \\
  Il \textbf{nostro} codice controlla il flusso di esecuzione.
  & È il \textbf{framework} a controllare il flusso di esecuzione, e invoca il
    nostro codice quando serve. \\
\end{tabular}
\end{center}

\subsection{Inversione del controllo}

\dfn{Inversion of Control}{
  L'\term{inversione del controllo} (IoC) è la caratteristica distintiva dei
  framework: il flusso di controllo è \textbf{invertito} rispetto alla
  programmazione imperativa tradizionale, e la responsabilità di gestirlo passa
  dallo sviluppatore al framework.}

È nota anche come \guillemotleft legge di Hollywood\guillemotright:
\textit{\guillemotleft Non chiamarci tu, ti chiamiamo noi\guillemotright}. La
figura~\ref{fig:ioc} segue una richiesta nei due casi, passo per passo.

\begin{figure}[H]
  \centering
  \begin{tikzpicture}[font=\Lato\scriptsize,
      b/.style={rounded corners=3pt, line width=0.8pt, align=center,
                minimum height=0.7cm},
      u/.style={b, draw=primary!65!black, fill=primary!10!white},
      l/.style={b, draw=neutral!55!white, fill=neutral!8!white, minimum width=1.5cm},
      f/.style={b, draw=orange-gradient-6, fill=orange-gradient-1!40!white},
      a/.style={-{Stealth[length=5pt,width=4pt]}, line width=0.8pt,
                draw=neutral!45!black},
      num/.style={circle, fill=primary!80!black, text=white,
                  font=\Lato\tiny\bfseries, inner sep=0pt, minimum size=10pt}]

    % con le librerie
    \node[font=\Lato\scriptsize\bfseries, text=neutral!25!black] at (2.0,2.55)
      {con le librerie: il nostro codice chiama};
    \node[u, minimum width=3.6cm] (uc) at (2.0,0) {codice utente};
    \node[l] (la) at (1.1,1.6) {libreria A};
    \node[l] (lb) at (2.9,1.6) {libreria B};
    \draw[a] (0.85,0.35) -- node[num, left=2pt]{2} (0.85,1.3);
    \draw[a, draw=accent!70!black] (1.35,1.3) -- node[num, right=2pt]{3} (1.35,0.35);
    \draw[a] (2.65,0.35) -- node[num, left=2pt]{4} (2.65,1.3);
    \draw[a, draw=accent!70!black] (3.15,1.3) -- node[num, right=2pt]{5} (3.15,0.35);
    \draw[a] (-1.0,0) -- node[num, above=2pt]{1} node[below=2pt]{richiesta} (uc.west);
    \draw[a] (uc.east) -- node[num, above=2pt]{6} node[below=2pt]{risposta} (5.0,0);

    % con un framework
    \node[font=\Lato\scriptsize\bfseries, text=neutral!25!black] at (9.5,2.55)
      {con un framework: il framework chiama il nostro codice};
    \node[f, minimum width=3.6cm] (fw) at (9.5,1.6) {\textbf{framework}};
    \node[u, minimum width=1.5cm] (u1) at (8.6,0) {codice\\utente};
    \node[u, minimum width=1.5cm] (u2) at (10.4,0) {codice\\utente};
    \draw[a] (8.35,1.25) -- node[num, left=2pt]{2} (8.35,0.4);
    \draw[a, draw=accent!70!black] (8.85,0.4) -- node[num, right=2pt]{3} (8.85,1.25);
    \draw[a] (10.15,1.25) -- node[num, left=2pt]{4} (10.15,0.4);
    \draw[a, draw=accent!70!black] (10.65,0.4) -- node[num, right=2pt]{5} (10.65,1.25);
    \draw[a] (6.5,1.6) -- node[num, above=2pt]{1} node[below=2pt]{richiesta} (fw.west);
    \draw[a] (fw.east) -- node[num, above=2pt]{6} node[below=2pt]{risposta} (12.5,1.6);
  \end{tikzpicture}
  \caption{A sinistra il nostro codice riceve la richiesta e usa le librerie,
           che gli restituiscono il controllo; a destra è il framework a
           ricevere la richiesta e a chiamare, quando serve, i pezzi di codice
           scritti da noi.}
  \label{fig:ioc}
\end{figure}

\subsection{Che cosa contiene un web framework}

\begin{itemize}
  \item \textbf{Funzionalità di base}: routing, analisi delle richieste,
        validazione degli input, gestione dei cookie, sessioni.
  \item \textbf{Motore di template}: aiuta a generare contenuto dinamico,
        separando la logica dalla presentazione.
  \item \textbf{Meccanismo ORM} (\textit{Object-Relational Mapping}):
        semplifica l'interazione con il database permettendo di lavorare con
        oggetti invece che con interrogazioni SQL.
\end{itemize}

\subsection{Framework opinionated}

I framework possono essere più o meno \term{opinionated}.

Un framework \textbf{fortemente opinionated} arriva con un insieme rigido di
convenzioni, buone pratiche e decisioni prese dai suoi autori: impone un modo
specifico di fare le cose e lascia poco margine allo sviluppatore. Un framework
\textbf{unopinionated} non impone un modo specifico.

\begin{center}
\renewcommand{\arraystretch}{1.4}
\begin{tabular}{@{}>{\raggedright\arraybackslash}p{6.3cm}@{\hspace{1.0cm}}>{\raggedright\arraybackslash}p{6.3cm}@{}}
  \textbf{Pro} & \textbf{Contro} \\
  Garantiscono maggiore coerenza fra progetti diversi.
  & Possono avere una curva di apprendimento più ripida. \\
  Accelerano lo sviluppo, riducendo l'affaticamento decisionale.
  & Possono non essere abbastanza flessibili per un certo progetto. \\
\end{tabular}
\end{center}

Di framework web lato server ce ne sono moltissimi, per ogni linguaggio:
Django e Flask (Python), Spring (Java), Laravel e Symfony (PHP), Ruby on Rails,
ASP.NET Core (C\#), e per Node.js Express, Koa, Hono, Fastify e altri ancora.

\section{Express}

\term{Express} si autodefinisce un \guillemotleft framework web veloce,
unopinionated e minimalista per Node.js\guillemotright, ed è di gran lunga il
framework di back-end più diffuso per Node.js (\textit{State of JavaScript
2024}). Si installa con npm, nella cartella del progetto:

\begin{console}[Terminale]
$ npm install express

added 65 packages, and audited 66 packages in 586ms

30 packages are looking for funding
  run `npm fund` for details

found 0 vulnerabilities
\end{console}

La versione installata oggi è \textbf{Express 5}: tutti gli esempi del
capitolo sono stati eseguiti con Express 5.3.0 e Node.js 22, e gli output
riportati sono quelli reali. (Il file \code{package.json} contiene
\code{"type": "module"}, per poter usare \code{import}, come nel
capitolo 11.)

\subsection{Il primo programma}

\begin{affianco}[0.55][c]
\begin{lstlisting}[style=js, name=index.js]
import express from "express";

const app = express();     // crea un'applicazione express
const PORT = 3000;

app.get('/', (req, res) => {
  res.send('Hello World!')
});

app.listen(PORT);
\end{lstlisting}
\risultato
\begin{browser}[localhost:3000]
  Hello World!
\end{browser}
\end{affianco}

Vale la pena confrontare queste otto righe con le decine di righe di
\code{switch} del capitolo 12. E Express fa più di quanto sembri: guardando la
risposta completa con \code{curl -i} (capitolo 9) si vedono le intestazioni
che ha aggiunto da solo, e chiedendo un percorso che non esiste si ottiene una
risposta \code{404} già pronta.

\begin{console}[Terminale]
$ curl -i http://localhost:3000/
HTTP/1.1 200 OK
X-Powered-By: Express
Content-Type: text/html; charset=utf-8
Content-Length: 12
ETag: W/"c-Lve95gjOVATpfV8EL5X4nxwjKHE"
Date: Tue, 03 Nov 2026 11:31:45 GMT
Connection: keep-alive
Keep-Alive: timeout=5

Hello World!
$ curl http://localhost:3000/nonesiste
<!DOCTYPE html>
<html lang="en">
<head>
<meta charset="utf-8">
<title>Error</title>
</head>
<body>
<pre>Cannot GET /nonesiste</pre>
</body>
</html>
\end{console}

\code{res.send} ha calcolato il tipo del contenuto (\code{text/html}, perché
gli abbiamo passato una stringa), la lunghezza e un \code{ETag}, l'impronta del
contenuto che il browser userà per la cache (capitolo 2). L'intestazione
\code{X-Powered-By} annuncia che il server usa Express: si può disattivare con
\code{app.disable("x-powered-by")}, per non dare informazioni gratuite a chi
cerca server vulnerabili.

\section{Routing}

Il \term{routing} è il modo in cui il framework stabilisce come gestire una
richiesta verso un dato \textit{endpoint}. Una rotta si definisce con i metodi
appropriati sull'oggetto \code{app}, tipicamente nella forma:

\begin{center}
  \code{app.<METODO>(<PERCORSO>, callback)}
\end{center}

dove \code{<METODO>} e \code{<PERCORSO>} identificano l'endpoint, e
\code{callback} è la funzione da eseguire per gestire le richieste che vi
arrivano.

\begin{affianco}[0.5][c]
\begin{lstlisting}[style=js, name=routing.js]
app.get("/hello", (req, res) => {
  res.send("Hello");
});

app.post("/hello", (req, res) => {
  res.send("Hello Post");
});
\end{lstlisting}
\risultato
\begin{immagine}[Stesso percorso, metodi diversi, risposte diverse.]
\begin{tikzpicture}[
    rq/.style={mcod, anchor=east, draw=primary!55!black, fill=primary!8!white,
      rounded corners=2pt, inner sep=3pt},
    rs/.style={mcod, anchor=west},
    a/.style={-{Stealth[length=4pt,width=3.5pt]}, line width=0.7pt, draw=neutral!45!black}]
  \node[rq] (r1) at (2.2,0) {GET /hello};
  \node[rs] (s1) at (3.0,0) {"Hello"};
  \node[rq] (r2) at (2.2,-0.65) {POST /hello};
  \node[rs] (s2) at (3.0,-0.65) {"Hello Post"};
  \node[rq] (r3) at (2.2,-1.3) {DELETE /hello};
  \node[rs] (s3) at (3.0,-1.3) {404};
  \draw[a] (r1) -- (s1); \draw[a] (r2) -- (s2); \draw[a, draw=merr] (r3) -- (s3);
\end{tikzpicture}
\end{immagine}
\end{affianco}

\begin{console}[Terminale]
$ curl http://localhost:3000/hello
Hello
$ curl -X POST http://localhost:3000/hello
Hello Post
$ curl -X DELETE http://localhost:3000/hello
...
<pre>Cannot DELETE /hello</pre>
...
\end{console}

L'opzione \code{-X} di \code{curl} sceglie il metodo della richiesta. Per un
metodo senza rotta Express risponde \code{404}, come per un percorso che non
esiste.

\subsection{Percorsi delle rotte}

I percorsi possono essere stringhe, schemi di stringa o espressioni regolari
complete. In caso di più corrispondenze, si applica \textbf{la prima} rotta che
corrisponde, nell'ordine in cui sono state dichiarate.

\begin{affianco}[0.58][c]
\begin{lstlisting}[style=js, name=paths.js]
app.get("/webtech", (req, res) => {
  res.send("Web Technologies!");
});

// corrisponde a /webtech, /wirelesstech,
// /wtech, /whatever-tech, ...
app.get(/^\/w.*tech$/, (req, res) => {
  res.send("Web Technologies Regex!");
});
\end{lstlisting}
\risultato
\begin{console}[Terminale]
$ curl localhost:3000/webtech
Web Technologies!
$ curl localhost:3000/wirelesstech
Web Technologies Regex!
$ curl localhost:3000/wtech
Web Technologies Regex!
$ curl localhost:3000/whatever-tech
Web Technologies Regex!
$ curl localhost:3000/webtechs
...
<pre>Cannot GET /webtechs</pre>
...
\end{console}
\end{affianco}

\code{/webtech} corrisponde a entrambe le rotte, ma vince la prima. L'espressione
regolare \code{/\^{}\textbackslash/w.*tech\$/} vuole una \code{w} subito dopo la
barra, poi qualunque sequenza di caratteri (\code{.*}) e infine \code{tech} in
fondo (\code{\$}): per questo \code{/webtechs} non corrisponde a niente.

\Needspace{14\baselineskip}
\begin{warningbox}{{Express 4 ed Express 5: gli schemi di stringa}}
  Molti esempi che si trovano in rete sono scritti per Express 4, che ammetteva
  nei percorsi i caratteri speciali \code{?}, \code{+}, \code{*} e le
  parentesi, per esempio \code{"/web(-)?tech*"}. In Express 5 non sono più
  validi, e il programma si ferma già all'avvio, quando la rotta viene
  dichiarata:
\begin{console}[Terminale]
$ node express4.js
...
PathError [TypeError]: Unexpected ( at index 4: /web(-)?tech*; visit https://git.new/pathToRegexpError for info
\end{console}
  La stessa idea, in Express 5, si scrive con le \textbf{graffe} per le parti
  facoltative e con un \textbf{jolly con nome}, \code{*nome}, per una sequenza
  qualsiasi:
\begin{lstlisting}[style=js, name=express5.js]
// {-} è una parte facoltativa, {*resto} un jolly facoltativo
app.get("/web{-}tech{*resto}", (req, res) => {
  res.send(`resto = ${JSON.stringify(req.params.resto)}`);
});
\end{lstlisting}
  Così \code{/webtech} e \code{/web-tech} corrispondono con \code{resto}
  indefinito, \code{/webtechnologies} con \code{resto} uguale a
  \code{["nologies"]}, e \code{/web-tech/corso/2026} con
  \code{["","corso","2026"]}: il jolly raccoglie i segmenti del percorso in un
  array. Chi vuole la massima libertà usa comunque un'espressione regolare.
\end{warningbox}

\subsection{Parametri di rotta}

Le rotte possono catturare segmenti di URL con un nome, detti \term{parametri},
dichiarati nel percorso con la sintassi \code{:nomeParametro}. I parametri
catturati sono resi disponibili nell'oggetto \code{req.params}.

\begin{affianco}[0.6][c]
\codicepiccolo
\begin{lstlisting}[style=js, name=params.js]
// per una richiesta a /student/42/exams/TecWeb
app.get("/student/:student_id/exams/:exam_name", (req, res) => {
  res.send(req.params);   // {"student_id":"42","exam_name":"TecWeb"}
});
\end{lstlisting}
\risultato
\begin{immagine}[I segmenti con i due punti diventano\\ proprietà di \code{req.params}.]
\begin{tikzpicture}[
    s/.style={mcod, anchor=base west, inner sep=1pt},
    p/.style={s, fill=accent!14!white}]
  \node[s] (a) at (0,0) {/student/};
  \node[p] (b) at (a.base east) {42};
  \node[s] (c) at (b.base east) {/exams/};
  \node[p] (e) at (c.base east) {TecWeb};
  \draw[wtbrace] ([yshift=-2pt]b.south west) -- ([yshift=-2pt]b.south east)
    node[midway, below=4pt, mcod]{student\_id};
  \draw[wtbrace] ([yshift=-2pt]e.south west) -- ([yshift=-2pt]e.south east)
    node[midway, below=4pt, mcod]{exam\_name};
  \mchiave{1.6cm}
  \oggetto{rp}{(1.7,-1.5)}{3.2cm}{req.params}{student\_id/"42", exam\_name/"TecWeb"}
\end{tikzpicture}
\end{immagine}
\end{affianco}

\begin{console}[Terminale]
$ curl -i http://localhost:3000/student/42/exams/TecWeb
HTTP/1.1 200 OK
X-Powered-By: Express
Content-Type: application/json; charset=utf-8
Content-Length: 40
...

{"student_id":"42","exam_name":"TecWeb"}
$ curl http://localhost:3000/student/bob/exams/TecWeb
{"student_id":"bob","exam_name":"TecWeb"}
\end{console}

Passando a \code{res.send} un oggetto, Express lo converte in JSON e imposta da
solo \code{Content-Type: application/json}. I parametri sono sempre
\textbf{stringhe}, e un parametro accetta qualunque segmento: anche \code{bob}
come matricola. Express 4 permetteva di vincolarne il formato con
un'espressione regolare fra parentesi (\code{:student\_id([0-9]+)}); Express 5
rifiuta anche questa sintassi, con lo stesso errore di prima. Si può
controllare il valore dentro il gestore, oppure scrivere tutta la rotta come
espressione regolare: i \textbf{gruppi con nome} (\code{(?<nome>\ldots)})
finiscono in \code{req.params}.

\begin{lstlisting}[style=js, name=params5.js]
app.get(/^\/student\/(?<student_id>[0-9]+)\/exams\/(?<exam_name>[^/]+)$/,
  (req, res) => {
    res.send(req.params);
  });
\end{lstlisting}

\begin{console}[Terminale]
$ curl localhost:3000/student/42/exams/TecWeb
{"student_id":"42","exam_name":"TecWeb"}
$ curl -o /dev/null -w '%{http_code}' localhost:3000/student/bob/exams/TecWeb
404
\end{console}

Le opzioni \code{-o /dev/null -w '\%\{http\_code\}'} dicono a \code{curl} di
buttare via il corpo e di stampare solo il codice di stato: \code{bob} non
corrisponde a \code{[0-9]+}, quindi la rotta non si applica e la risposta è
\code{404}.

\info{Il problema risolto}{
  Questo è esattamente il caso \code{GET /users/42/friends} che nel capitolo 12
  avevamo indicato come impossibile da gestire con uno \code{switch} su stringhe
  esatte. Con Express diventa una riga.}

\subsection{Più gestori per una rotta}

È possibile specificare più funzioni gestore per una rotta. Si comportano come
\textit{middleware} (li vedremo fra poco), e \code{next} è la callback che
invoca il gestore successivo.

\begin{affianco}[0.6][t]
\codicepiccolo
\begin{lstlisting}[style=js, name=handlers.js]
let f1 = function (req, res, next) { console.log("f1"); next(); }
let f2 = function (req, res, next) { console.log("f2"); next(); }
let f3 = function (req, res, next) { console.log("f3"); next(); }

app.get("/handlers", [f1, f2, f3], (req, res) => {
  res.send("Done with all handlers");
});
\end{lstlisting}
\begin{immagine}
\begin{tikzpicture}[
    f/.style={mfun},
    a/.style={-{Stealth[length=4pt,width=3.5pt]}, line width=0.7pt, draw=accent!70!black}]
  \node[f] (f1) at (0,0) {f1};
  \node[f] (f2) at (1.6,0) {f2};
  \node[f] (f3) at (3.2,0) {f3};
  \node[f, fill=green-gradient-1!45!white, draw=green-gradient-6] (h) at (5.1,0) {gestore};
  \draw[a] (f1) -- node[mcod, above, font=\FiraCodeNL\tiny]{next()} (f2);
  \draw[a] (f2) -- node[mcod, above, font=\FiraCodeNL\tiny]{next()} (f3);
  \draw[a] (f3) -- node[mcod, above, font=\FiraCodeNL\tiny]{next()} (h);
\end{tikzpicture}
\end{immagine}
\risultato
\begin{console}[Terminale]
$ curl localhost:3000/handlers
Done with all handlers
\end{console}
\begin{console}[Terminale del server]
f1
f2
f3
\end{console}
\end{affianco}

Il client vede solo la risposta dell'ultimo gestore; le tre righe di
\code{console.log} compaiono nel terminale del server.

\subsection{Rotte concatenabili e router modulari}

Con \code{app.route()} il percorso si specifica in un solo punto (meno
ridondanza) e i gestori specifici per metodo si concatenano.

\begin{affianco}[0.5][c]
\begin{lstlisting}[style=js, name=book.js]
app.route('/book')
  .get((req, res) => {
    res.send('Get a random book')
  })
  .post((req, res) => {
    res.send('Add a book')
  })
  .delete((req, res) => {
    res.send('Delete a book')
  });
\end{lstlisting}
\risultato
\begin{console}[Terminale]
$ curl localhost:3000/book
Get a random book
$ curl -X POST localhost:3000/book
Add a book
$ curl -X DELETE localhost:3000/book
Delete a book
$ curl -X PUT localhost:3000/book
...
<pre>Cannot PUT /book</pre>
...
\end{console}
\end{affianco}

Si può anche creare un \term{router modulare} (a sinistra), da montare
nell'applicazione principale (a destra):

\begin{affianco}[0.45]
\begin{lstlisting}[style=js, name=router.js]
import express from "express";

export const router = express.Router();

router.get("/echo/:value", (req, res) => {
  res.send(req.params.value);
});
\end{lstlisting}
\begin{console}[Terminale]
$ curl localhost:3000/custom/echo/ciao
ciao
$ curl localhost:3000/echo/ciao
...
<pre>Cannot GET /echo/ciao</pre>
...
\end{console}
\risultato
\begin{lstlisting}[style=js, name=index.js]
import express from "express";
import { router as echoRouter } from "./router.js";

const app = express();

// monta echoRouter, opzionalmente sotto un certo percorso
app.use("/custom", echoRouter);
// /custom/echo/:value è gestita dal router

app.listen(3000);
\end{lstlisting}
\end{affianco}

Il router è montato sotto \code{/custom}: i suoi percorsi sono
\textbf{relativi} a quel prefisso, quindi \code{/echo/ciao} da solo non esiste.
Così si tiene ordinata un'applicazione reale: un file di rotte per ciascuna
area funzionale, montato dove serve.

\section{Rispondere al client}

Sull'oggetto risposta (\code{res}) vanno invocati i metodi appropriati per
inviare una risposta al client, \textbf{altrimenti la richiesta resta appesa}.

\begin{center}
\renewcommand{\arraystretch}{1.35}
\begin{tabular}{@{}l@{\hspace{1.0cm}}>{\raggedright\arraybackslash}p{8.4cm}@{}}
  \textbf{Metodo} & \textbf{Descrizione} \\
  \code{res.download()}   & propone il download di un file \\
  \code{res.end()}        & termina il processo di risposta \\
  \code{res.json()}       & invia una risposta JSON \\
  \code{res.jsonp()}      & invia una risposta JSON con supporto JSONP \\
  \code{res.redirect()}   & reindirizza una richiesta \\
  \code{res.render()}     & genera un template di vista \\
  \code{res.send()}       & invia una risposta di vario tipo \\
  \code{res.sendFile()}   & invia un file come flusso di byte \\
  \code{res.sendStatus()} & imposta il codice di stato e ne invia la
                            rappresentazione testuale come corpo \\
\end{tabular}
\end{center}

\subsection{Generare un template}

Basta dire a Express quale motore di template usare: \code{res.render} cerca il
template nella cartella \code{views}, lo trasforma in HTML e lo invia.

\begin{lstlisting}[style=js, name=template.js]
import express from "express";

const app = express();

app.set("view engine", "pug");

app.get('/template', (req, res) => {
  res.render("template");
});
\end{lstlisting}

\begin{affianco}[0.55][c]
\begin{lstlisting}[style=pug, name=views/template.pug]
doctype html
html
  head
    title Pug Template rendered in Express
  body
    h1 Rendering a Pug Template in Express
    p Il procedimento è immediato.
\end{lstlisting}
\risultato
\begin{browser}[Pug Template rendered in Express]
  {\raggedright\Lora\bfseries\LARGE Rendering a Pug Template in Express\par}\vspace{6pt}
  Il procedimento è immediato.
\end{browser}
\end{affianco}

\subsection{Passare dati alla vista}

I dati impostati nell'oggetto \code{res.locals} sono disponibili a tutte le
viste generate per quella risposta.

\begin{lstlisting}[style=js, name=template.js, firstnumber=10]

app.get("/example", (req, res) => {
  res.locals.items = [
    { name: "Web tech.", grade: "A+" },
    { name: "Soft. Eng.", grade: "A" },
    { name: "Calculus",  grade: "B" },
  ];
  res.render("example");
});

app.listen(3000);
\end{lstlisting}

\begin{affianco}[0.5][c]
\begin{lstlisting}[style=pug, name=views/example.pug]
doctype html
html
  head
    title List of Exams
  body
    h1 List of Exams
    each exam in items
      li #{exam.name} grade #{exam.grade}
\end{lstlisting}
\risultato
\begin{browser}[List of Exams]
  {\Lora\bfseries\LARGE List of Exams}
  \begin{itemize}[leftmargin=1.8em, itemsep=0pt, topsep=4pt, parsep=0pt]
    \item Web tech. grade A+
    \item Soft. Eng. grade A
    \item Calculus grade B
  \end{itemize}
\end{browser}
\end{affianco}

Al browser arriva l'HTML generato da Pug, tutto su una riga:

\begin{console}[Terminale]
$ curl localhost:3000/example
<!DOCTYPE html><html><head><title>List of Exams</title></head><body><h1>List of Exams</h1><li>Web tech. grade A+</li><li>Soft. Eng. grade A</li><li>Calculus grade B</li></body></html>
\end{console}

Gli elementi \code{li} non sono dentro una lista \code{ul}, perché il template
delle slide non la mette: il browser li mostra comunque come voci di elenco, ma
un \code{ul} attorno al ciclo \code{each} renderebbe l'HTML corretto.

\section{Il ciclo richiesta-risposta}

Il ciclo richiesta-risposta è il cuore di Express. Quando arriva una richiesta:

\begin{enumerate}
  \item vengono creati un oggetto richiesta e un oggetto risposta;
  \item viene eseguita una \textbf{sequenza} di uno o più passi di elaborazione
        (analizzare il corpo, analizzare i cookie, stabilire una sessione,
        \ldots);
  \item viene preparata e inviata la risposta.
\end{enumerate}

\begin{figure}[H]
  \centering
  \begin{tikzpicture}[font=\Lato\scriptsize,
      m/.style={rounded corners=3pt, draw=primary!65!black, line width=0.8pt,
                fill=primary!10!white, minimum width=2.3cm, minimum height=1.5cm,
                align=center},
      h/.style={m, draw=accent!60!black, fill=accent!12!white},
      a/.style={-{Stealth[length=5pt,width=4pt]}, line width=0.85pt,
                draw=neutral!45!black},
      r/.style={-{Stealth[length=5pt,width=4pt]}, line width=0.85pt,
                draw=accent!70!black}]

    \node[font=\Lato\scriptsize, anchor=east, align=right, text=neutral!25!black]
      (in) at (-0.5,0.4) {richiesta\\ HTTP};

    \node[m] (m1) at (1.4,0.4) {Middleware\\[2pt]\scriptsize log della\\\scriptsize richiesta};
    \node[m] (m2) at (4.2,0.4) {Middleware\\[2pt]\scriptsize analisi del\\\scriptsize corpo};
    \node[m] (m3) at (7.0,0.4) {Middleware\\[2pt]\scriptsize controllo\\\scriptsize autenticazione};
    \node[h] (rh) at (9.8,0.4) {Gestore\\ della rotta\\[2pt]\scriptsize \code{res.render()}};

    \draw[a] (in.east) -- (m1.west);
    \draw[a] (m1) -- node[wtlab, above]{\code{next()}} (m2);
    \draw[a] (m2) -- node[wtlab, above]{\code{next()}} (m3);
    \draw[a] (m3) -- node[wtlab, above]{\code{next()}} (rh);

    \node[font=\Lato\scriptsize, anchor=west, align=left, text=accent!60!black]
      (out) at (11.3,-1.5) {risposta\\ HTTP};
    \foreach \n in {m1,m2,m3,rh} {
      \draw[r] (\n.south) -- ++(0,-0.55);
    }
    \draw[r, draw=accent!70!black] (1.4,-0.9) -- (10.9,-0.9) -- (out.west);
    \node[font=\Lato\scriptsize\itshape, text=accent!55!black] at (5.5,-1.2)
      {ogni livello può rispondere e interrompere la catena};

    \begin{scope}[on background layer]
      \node[rounded corners=5pt, draw=primary!35!white, line width=0.9pt,
            fit=(m1)(m2)(m3), inner sep=9pt,
            label={[font=\Lato\scriptsize, text=primary!60!black]above:pila dei
            middleware}] {};
    \end{scope}
  \end{tikzpicture}
  \caption{Il ciclo richiesta-risposta di Express.}
  \label{fig:express-cycle}
\end{figure}

\section{I middleware}

\dfn{Middleware}{
  In Express i \term{middleware} sono funzioni che possono manipolare gli
  oggetti richiesta e risposta correnti, oppure eseguire codice arbitrario. Si
  chiamano così perché sono eseguiti \textbf{nel mezzo}, fra la ricezione di una
  richiesta e l'invio della relativa risposta.}

I middleware:

\begin{itemize}
  \item vengono eseguiti \textbf{nello stesso ordine in cui sono dichiarati};
  \item ricevono come argomenti \code{req} (l'oggetto richiesta), \code{res}
        (l'oggetto risposta) e \code{next} (una callback che invoca il
        middleware successivo nella pila).
\end{itemize}

\begin{esempio}{Il primo middleware}
\begin{affianco}[0.45]
\begin{lstlisting}[style=js, name=middleware.js]
export function logger(req, res, next) {
  console.log("LOGGER HAS BEEN CALLED!");
  next();
}
\end{lstlisting}
\begin{console}[Terminale]
$ curl localhost:3000/
Hello World!
$ curl localhost:3000/ciao
...
<pre>Cannot GET /ciao</pre>
...
\end{console}
\risultato
\begin{lstlisting}[style=js, name=index.js]
app.use(logger);   // monta logger sul percorso "/"

app.get('/', (req, res) => {
  res.send('Hello World!')
});

app.listen(PORT);
\end{lstlisting}
\begin{console}[Terminale del server]
LOGGER HAS BEEN CALLED!
LOGGER HAS BEEN CALLED!
\end{console}
\end{affianco}

Montato con \code{app.use} sul percorso \code{/}, il middleware viene eseguito
per \textbf{ogni} richiesta: anche per \code{/ciao}, che poi finisce in un
\code{404}. Un middleware non sa nulla delle rotte che vengono dopo di lui.
\end{esempio}

\warning{Dimenticare \code{next()}}{
  Se un middleware non invoca \code{next()} e non invia una risposta, la
  richiesta resta \textbf{appesa} finché il client non va in timeout. È l'errore
  più comune con Express, e non produce alcun messaggio d'errore: semplicemente
  la pagina non carica mai.}

\begin{esempio}{Un middleware che arricchisce la richiesta}
\begin{affianco}[0.52]
\begin{lstlisting}[style=js, name=middleware.js, firstnumber=5]

export function addTimestamp(req, res, next) {
  req.timestamp = new Date();
  next();
}
\end{lstlisting}
\begin{lstlisting}[style=js, name=index.js]
app.use(addTimestamp);
app.use(logger);

app.get('/', (req, res) => {
  res.send(`Richiesta ricevuta il ${req.timestamp}`)
});
\end{lstlisting}
\risultato
\begin{immagine}[\code{addTimestamp} arricchisce \code{req}.]
\begin{tikzpicture}
  \mchiave{1.4cm}
  \oggetto{r}{(0,0)}{3.4cm}{req}{method/"GET", url/"/", timestamp/Date \ldots}
  \node[draw=accent!60!black, line width=1.1pt, fit=(r-3), inner sep=0pt] {};
\end{tikzpicture}
\end{immagine}
\begin{console}[Terminale]
$ curl localhost:3000/
Richiesta ricevuta il Tue Nov 03 2026 12:32:56 GMT+0100 (Central European Standard Time)
\end{console}
\end{affianco}

Qui si vede il punto centrale: un middleware può \textbf{aggiungere proprietà}
all'oggetto richiesta, e tutti i middleware e i gestori successivi le trovano
già pronte. È il meccanismo con cui l'oggetto \code{context} che avevamo
passato a mano nel capitolo 13 diventa superfluo.
\end{esempio}

\section{Gestione degli errori}

\subsection{Errori nel codice sincrono}

È importante assicurarsi che ogni errore verificatosi durante l'elaborazione dei
middleware o delle rotte sia gestito correttamente da Express. Se un errore
viene sollevato durante l'esecuzione di codice \textbf{sincrono}, viene
catturato automaticamente dal gestore di errori predefinito di Express.

\begin{affianco}[0.5][c]
\begin{lstlisting}[style=js, name=errors.js]
import express from "express";
import fs from "fs";

const app = express();

app.get("/error", (req, res) => {
  throw new Error("Si è verificato un errore");
});
\end{lstlisting}
\risultato
\begin{browser}[Error]
  {\FiraCodeNL\scriptsize Error: Si è verificato un errore\\
   \hspace*{1em}at file:///[cartella del progetto]/errors.js:7:9\\
   \hspace*{1em}at Layer.handleRequest ([\ldots]/router/lib/layer.js:152:17)\\
   \hspace*{1em}at next ([\ldots]/router/lib/route.js:157:13)\\
   \hspace*{1em}\ldots}
\end{browser}
\wtnota{Stato 500; la traccia è quella reale, abbreviata.}
\end{affianco}

Il server non si ferma: risponde con lo stato \code{500} e una pagina con la
\textbf{traccia dello stack} (stampata anche nel suo terminale), che indica
dove è nato l'errore: \code{errors.js}, riga 7, colonna 9. La traccia compare
solo in \textbf{sviluppo}: in produzione (variabile d'ambiente
\code{NODE\_ENV} uguale a \code{production}) il client riceve un messaggio
generico, perché i dettagli interni non devono arrivare agli utenti.

\begin{console}[Terminale]
$ NODE_ENV=production node errors.js &
$ curl localhost:3000/error
...
<pre>Internal Server Error</pre>
...
\end{console}

\subsection{Errori nel codice asincrono}

Gli errori che si verificano nelle callback di funzioni \textbf{asincrone},
invocate da gestori di rotta e middleware, devono invece essere passati
esplicitamente a \code{next()}: quando la callback viene eseguita, il gestore
è già terminato, ed Express non può catturare ciò che viene sollevato.

\begin{affianco}[0.5][c]
\begin{lstlisting}[style=js, name=errors.js, firstnumber=9]

app.get('/readfile', (req, res, next) => {
  fs.readFile('/file-does-not-exist', (err, data) => {
    if (err) {
      next(err);       // passa l'errore a Express
    } else {
      res.send(data);
    }
  })
});
\end{lstlisting}
\risultato
\begin{immagine}[L'errore della callback va consegnato\\ a mano a \code{next}.]
\begin{tikzpicture}[
    st/.style={rounded corners=3pt, draw=primary!55!black, fill=primary!6!white,
      line width=0.7pt, font=\FiraCodeNL\scriptsize, inner sep=3pt},
    a/.style={-{Stealth[length=4pt,width=3.5pt]}, line width=0.7pt, draw=merr}]
  \node[st] (r) at (0,0) {fs.readFile(\ldots)};
  \node[st, draw=merr, fill=red-gradient-1!35!white] (n) at (0,-1.1) {next(err)};
  \node[st] (h) at (0,-2.2) {gestore degli errori};
  \draw[a] (r) -- node[mnota, right]{\code{err}} (n);
  \draw[a] (n) -- (h);
\end{tikzpicture}
\end{immagine}
\end{affianco}

\begin{console}[Terminale]
$ curl -i localhost:3000/readfile
HTTP/1.1 500 Internal Server Error
...
<pre>Error: ENOENT: no such file or directory, open &#39;/file-does-not-exist&#39;</pre>
...
\end{console}

\code{ENOENT} (\textit{Error NO ENTry}) è il codice con cui il sistema
operativo dice che il file non esiste; \code{\&\#39;} è l'apostrofo, scritto
come riferimento a carattere (capitolo 2) perché il gestore predefinito fa
l'escaping del messaggio.

\subsection{Gestori di errore personalizzati}

Express include un gestore di errori predefinito, quello visto sopra, che
stampa l'intera traccia dello stack solo in modalità sviluppo. Se ne possono
definire di propri (il programma della prossima sezione ne usa uno): sono
middleware speciali con \textbf{quattro} parametri di ingresso, e vanno
dichiarati \textbf{dopo} tutte le rotte.

\begin{lstlisting}[style=js]
app.use((err, req, res, next) => {
  console.error(err.stack);
  res.status(500).send('Something broke!');
});
\end{lstlisting}

\info{Quattro parametri, non tre}{
  È l'\textbf{arità} della funzione a dire a Express che si tratta di un gestore
  di errori: una funzione con tre parametri è un middleware normale, una con
  quattro è un gestore di errori. Ometterne uno per distrazione fa sì che il
  gestore non venga mai invocato, senza alcun avviso.}

\subsection{Le Promise in Express 5}

Da \textbf{Express 5} in poi, i gestori e i middleware che restituiscono una
Promise invocano automaticamente \code{next(value)} quando questa viene
rifiutata o solleva un errore: non serve farlo esplicitamente. Una funzione
\code{async} restituisce sempre una Promise (capitolo 8), quindi con
\code{await} si scrive codice asincrono come se fosse sincrono, e gli errori
arrivano comunque al gestore.

\begin{affianco}[0.57]
\codicepiccolo
\begin{lstlisting}[style=js, name=async5.js]
import express from "express";

const app = express();

const users = new Map([["1", { id: 1, name: "Janet" }]]);

// simula una lettura dal database, che richiede tempo
function getUserById(id) {
  return new Promise((resolve, reject) => {
    setTimeout(() => {
      if (users.has(id)) {
        resolve(users.get(id));
      } else {
        reject(new Error(`Utente ${id} non trovato`));
      }
    }, 100);
  });
}

app.get('/user/:id', async (req, res, next) => {
  const user = await getUserById(req.params.id)
  res.send(user)
});

app.use((err, req, res, next) => {
  console.error(err.stack);
  res.status(500).send('Something broke!');
});

app.listen(3000);
\end{lstlisting}
\risultato
\begin{console}[Terminale]
$ curl localhost:3000/user/1
{"id":1,"name":"Janet"}
$ curl -i localhost:3000/user/99
HTTP/1.1 500 Internal Server Error
...

Something broke!
\end{console}
\begin{console}[Terminale del server]
Error: Utente 99 non trovato
    at Timeout._onTimeout (file:///[cartella del progetto]/async5.js:14:16)
    at listOnTimeout (node:internal/timers:588:17)
    at process.processTimers (node:internal/timers:523:7)
\end{console}
\end{affianco}

Per l'utente 99 la Promise viene rifiutata, \code{await} solleva l'errore
dentro il gestore, ed Express lo consegna al \textbf{gestore di errori}
personalizzato in fondo al programma, che stampa la traccia nel terminale del
server e risponde \code{Something broke!}.

\subsection{Le Promise in Express 4}

Con Express 4 lo stesso programma si comporta in modo molto diverso, ed è bene
saperlo perché molto codice in circolazione è scritto per quella versione:
Express 4 ignora la Promise restituita dal gestore, il rifiuto non viene
gestito da nessuno, e Node.js, per un rifiuto non gestito, \textbf{termina
l'intero processo}. Nella prova (Express 4.22 e Node.js 22) la richiesta
per l'utente 99 ha fatto cadere il server:

\begin{console}[Terminale del server, con Express 4]
$ node async5.js
file:///[cartella del progetto]/async5.js:14
        reject(new Error(`Utente ${id} non trovato`));
               ^

Error: Utente 99 non trovato
    at Timeout._onTimeout (file:///[cartella del progetto]/async5.js:14:16)
    ...

Node.js v22.22.0
$
\end{console}

Con Express 4 il gestore va quindi scritto con \code{try}/\code{catch}, e
nel \code{catch} si chiama \code{next(err)} a mano.

\section{Middleware utili}

\subsection{express.static}

I primi tre middleware di questa sezione sono integrati in Express; nelle prove
stanno tutti nello stesso programma. \code{express.static} è progettato per
servire file statici. Usa \code{req.url}
per determinare se una richiesta corrisponde a un file della directory radice
indicata.

\begin{affianco}[0.55][c]
\begin{lstlisting}[style=js]
// i file della directory public diventano statici
app.use(express.static("public"));             // montato su "/"
app.use("/static", express.static("public"));  // sul percorso virtuale "/static"
\end{lstlisting}
\risultato
\begin{immagine}[Ogni file di \code{public} ha la sua URL.]
\begin{tikzpicture}[
    f/.style={mcod, anchor=west},
    u/.style={mcod, anchor=west, text=primary!60!black},
    a/.style={-{Stealth[length=4pt,width=3.5pt]}, line width=0.7pt, draw=neutral!45!black}]
  \node[f] at (0,0) {public/};
  \foreach \p/\u [count=\i] in {css\slash style.css/\slash css\slash style.css,
      index.html/\slash{} oppure \slash index.html} {
    \node[f] (p\i) at (0.4,-\i*0.5) {\p};
    \node[u] (u\i) at (3.0,-\i*0.5) {\u};
    \draw[a] (p\i.east -| 2.5,0) -- (u\i.west);
  }
\end{tikzpicture}
\end{immagine}
\end{affianco}

\begin{console}[Terminale]
$ curl -i localhost:3000/css/style.css
HTTP/1.1 200 OK
X-Powered-By: Express
Accept-Ranges: bytes
Cache-Control: public, max-age=0
Last-Modified: Tue, 03 Nov 2026 11:33:23 GMT
ETag: W/"22-1a1a12f71d3"
Content-Type: text/css; charset=utf-8
Content-Length: 34
...

body { font-family: sans-serif; }
$ curl localhost:3000/static/css/style.css
body { font-family: sans-serif; }
$ curl localhost:3000/
<!DOCTYPE html>
<title>Benvenuti</title>
<h1>Pagina statica</h1>
\end{console}

Il middleware ricava il \code{Content-Type} dall'estensione del file, e aggiunge
le intestazioni per la cache (\code{Last-Modified}, \code{ETag}). Per il
percorso \code{/} serve il file \code{index.html} della cartella, come fanno i
server web. Sostituisce da solo tutti i \code{case} con \code{fs.readFile} che
avevamo scritto a mano nel capitolo 12.

\subsection{express.urlencoded}

Middleware integrato che analizza il corpo delle richieste con contenuto
codificato come URL, cioè quelle inviate dai form, e ne mette i parametri
nell'oggetto \code{req.body}.

\begin{affianco}[0.55]
\begin{lstlisting}[style=js]
app.use(express.urlencoded({ extended: false }));

app.get("/form", (req, res) => { res.render("form"); });

app.post("/form", (req, res) => {
  res.send(`msg=${req.body.msg}; num=${req.body.num}`);
});
\end{lstlisting}
\risultato
\begin{lstlisting}[style=pug, name=views/form.pug]
h1 Form
form(action="/form" method="POST")
  input(type="text" name="msg")
  input(type="number" name="num")
  input(type="submit")
\end{lstlisting}
\begin{browser}[localhost:3000/form]
  {\Lora\bfseries\LARGE Form}\\[4pt]
  \wtcampo[2.0cm]{Hello}\ \wtcampo[1.0cm]{42}\ \wtbottone{Invia}
\end{browser}
\end{affianco}

\begin{console}[Terminale]
$ curl -d 'msg=Hello&num=42' localhost:3000/form
msg=Hello; num=42
\end{console}

Con \code{-d} \code{curl} invia una \code{POST} con il corpo indicato, codificato
come quello di un form: è ciò che fa il browser premendo il pulsante. Una riga
sostituisce le due funzioni di analisi del corpo scritte a mano nel
capitolo 12, gestione della codifica URL compresa. L'opzione
\code{extended: false} sceglie il formato semplice, \code{chiave=valore}; con
\code{true} si potrebbero inviare anche oggetti annidati (\code{a[b]=1}).

\subsection{express.json}

Si comporta come \code{express.urlencoded}, ma analizza corpi in formato JSON.
Sarà fondamentale nel capitolo sulle API REST.

\begin{lstlisting}[style=js]
app.use(express.json());

app.post("/json", (req, res) => {
  res.send(`exam=${req.body.exam}; grade=${req.body.grade}`);
});
\end{lstlisting}

\begin{affianco}[0.5]
\begin{lstlisting}[style=http, name=Richiesta]
POST http://localhost:3000/json HTTP/1.1
Content-Type: application/json

{
  "exam": "Web Technologies",
  "grade": "A+"
}
\end{lstlisting}
\risultato
\begin{lstlisting}[style=http, name=Risposta]
HTTP/1.1 200 OK
X-Powered-By: Express
Content-Type: text/html; charset=utf-8
Content-Length: 31
ETag: W/"1f-hq1l5pZKHQLA+ZKKA/JO5XfvqTE"
Date: Tue, 03 Nov 2026 11:33:23 GMT
Connection: keep-alive
Keep-Alive: timeout=5

exam=Web Technologies; grade=A+
\end{lstlisting}
\end{affianco}

La richiesta è stata inviata con
\code{curl -i -H 'Content-Type: application/json' -d '\{"exam": \ldots\}'}:
l'opzione \code{-H} aggiunge l'intestazione. È proprio il \code{Content-Type}
a decidere quale dei due middleware analizza il corpo: ciascuno ignora le
richieste che non sono nel suo formato.

\subsection{express-session}

Middleware disponibile su npm (\code{npm install express-session}), semplifica
la gestione di una sessione memorizzata sul server: i dati restano sul server,
l'identificatore viaggia in un cookie. È esattamente il meccanismo che nel
capitolo 13 abbiamo implementato a mano.

\begin{affianco}[0.6][c]
\codicepiccolo
\begin{lstlisting}[style=js, name=session.js]
import express from "express";
import session from "express-session";

const app = express();

app.use(session({
  secret: "s3cr37",
  saveUninitialized: false,
  resave: false
}));

app.get("/session-counter", (req, res) => {
  if (!req.session?.count) {
    req.session.count = 1;
    res.send("È la prima volta che visiti questa pagina.");
  } else {
    req.session.count = req.session.count + 1;
    res.send(`Hai visitato questa pagina ${req.session.count} volte.`);
  }
});
\end{lstlisting}
\risultato
\begin{immagine}[Il contatore sta sul server; il browser\\ conserva solo il cookie \code{connect.sid}.]
\begin{tikzpicture}
  \mchiave{0.9cm}
  \oggetto[fill=green-gradient-1!50!white]{s}{(0,0)}{2.4cm}{sessione sul server}{count/3}
  \node[mcod, anchor=east] (c) at (-1.9,-0.25) {connect.sid=s\%3A\ldots};
  \draw[-{Stealth[length=4pt,width=3.5pt]}, line width=0.7pt, draw=accent!70!black] (c.east) -- (s-box.west |- c);
\end{tikzpicture}
\end{immagine}
\end{affianco}

Con \code{curl} il browser si simula così: \code{-c jar} salva i cookie
ricevuti nel file \code{jar}, \code{-b jar} li rimanda con la richiesta.

\begin{console}[Terminale]
$ curl -i -c jar localhost:3000/session-counter
HTTP/1.1 200 OK
...
Set-Cookie: connect.sid=s%3ADkHKSsnn4-w7TNs24QCDH4nC37QuuMc1.JJ22WTJ%2Bb8SCxGjjNeFqJ6SIS8scb3cEBI92vGMuM%2Bk; Path=/; HttpOnly
...

È la prima volta che visiti questa pagina.
$ curl -b jar -c jar localhost:3000/session-counter
Hai visitato questa pagina 2 volte.
$ curl -b jar -c jar localhost:3000/session-counter
Hai visitato questa pagina 3 volte.
$ curl localhost:3000/session-counter
È la prima volta che visiti questa pagina.
\end{console}

\begin{itemize}
  \item La quarta richiesta non manda il cookie: per il server è un altro
        client, con una sessione nuova.
  \item Il valore del cookie, decodificato, è \code{s:}, l'identificatore
        della sessione (\code{DkHKSs\ldots}), un punto e una \textbf{firma}
        calcolata con il \code{secret}. Se il client modifica anche un solo
        carattere dell'identificatore, la firma non corrisponde più e il server
        lo ignora: nella prova, cambiando un carattere del cookie, la risposta è
        tornata \guillemotleft È la prima volta\ldots\guillemotright. Per
        questo il \code{secret} non va mai reso pubblico.
  \item \code{saveUninitialized: false} non crea sessioni (né cookie) per chi
        non vi ha ancora scritto nulla; \code{resave: false} non riscrive la
        sessione nell'archivio se la richiesta non l'ha modificata.
\end{itemize}

\warning{L'archivio predefinito non è per la produzione}{
  L'archivio di sessione predefinito di \code{express-session} tiene tutto in
  memoria: per la produzione va configurato un archivio dedicato, per esempio
  basato su \textit{memcached} o \textit{Redis}. Sono gli stessi limiti già
  discussi nel capitolo 13.}

\subsection{cookie-parser}

Altro middleware disponibile su npm: analizza l'header \code{Cookie} delle
richieste in arrivo e popola \code{req.cookies} con un oggetto le cui proprietà
sono i nomi dei cookie.

\begin{affianco}[0.6][c]
\codicepiccolo
\begin{lstlisting}[style=js, name=cookies.js]
import express from "express";
import cookieParser from "cookie-parser";

const app = express();

app.use(cookieParser());

app.get("/cookie-counter", (req, res) => {
  if (!req.cookies?.count) {
    res.cookie("count", 1, { maxAge: 10 * 60 * 1000 });
    res.send("È la prima volta che visiti questa pagina.");
  } else {
    let count = parseInt(req.cookies.count) + 1;
    res.cookie("count", count, { maxAge: 10 * 60 * 1000 });
    res.send(`Hai visitato questa pagina ${count} volte.`);
  }
});
\end{lstlisting}
\risultato
\begin{immagine}[Qui invece il contatore viaggia nel cookie:\\
                 il client può leggerlo e cambiarlo.]
\begin{tikzpicture}
  \mchiave{1.0cm}
  \oggetto[fill=neutral!12!white]{c}{(0,0)}{3.0cm}{cookie nel browser}{count/"3"}
\end{tikzpicture}
\end{immagine}
\end{affianco}

\begin{console}[Terminale]
$ curl -i -c jar2 localhost:3000/cookie-counter
HTTP/1.1 200 OK
X-Powered-By: Express
Set-Cookie: count=1; Max-Age=600; Path=/; Expires=Tue, 03 Nov 2026 11:43:32 GMT
...

È la prima volta che visiti questa pagina.
$ curl -b jar2 -c jar2 localhost:3000/cookie-counter
Hai visitato questa pagina 2 volte.
$ curl -b jar2 -c jar2 localhost:3000/cookie-counter
Hai visitato questa pagina 3 volte.
$ curl -b 'count=41' localhost:3000/cookie-counter
Hai visitato questa pagina 42 volte.
\end{console}

\code{res.cookie} imposta il cookie con la sua durata: \code{maxAge} è in
millisecondi, \code{Max-Age} nell'intestazione in secondi (dieci minuti). L'ultima
richiesta mostra la differenza con le sessioni: il valore sta nel client, e
basta inviarne uno inventato (\code{-b 'count=41'}) perché il server ci creda.

\section{Messaggi flash}

Spesso, nelle applicazioni web, gli utenti vengono \textbf{reindirizzati} dopo
aver compiuto un'azione, per esempio riportati alla homepage dopo aver
cancellato dei dati. È buona norma fornire un riscontro, così che sappiano se
l'operazione è andata a buon fine.

La parte difficile è che, dopo il reindirizzamento, comincia un ciclo
richiesta-risposta \textbf{completamente nuovo}. Come tenere traccia del fatto
che dobbiamo mostrare un messaggio?

È l'ennesima variazione del solito problema: HTTP non ha stato. Si risolve nei
modi già visti (cookie, sessioni), e alcuni middleware sono molto utili allo
scopo. Uno è \term{connect-flash}, disponibile su npm: salva i messaggi
\textbf{nella sessione}, quindi va registrato dopo \code{express-session}.

\begin{lstlisting}[style=js, name=flash.js]
import express from "express";
import session from "express-session";
import flash from "connect-flash";

const app = express();
app.set("view engine", "pug");

app.use(session({ secret: "s3cr37", saveUninitialized: false, resave: false }));
app.use(flash());

app.get("/redirect", (req, res) => {
  req.flash("info", "Sei stato reindirizzato.");
  req.flash("info", "Questa è una dimostrazione dei messaggi flash.");
  req.flash("success", "È andato tutto bene.");
  req.flash("error", "Ma abbiamo anche un messaggio di errore.");
  res.redirect("/showflash");
});

app.get("/showflash", (req, res) => {
  let infos     = req.flash("info");      // consuma i messaggi con quella chiave
  let successes = req.flash("success");
  let errors    = req.flash("error");
  res.render("flash", { infos: infos, successes: successes, errors: errors });
});
\end{lstlisting}

\begin{affianco}[0.5][c]
\begin{lstlisting}[style=pug, name=views/flash.pug]
ul
  each info in infos
    li(style="color: blue") #{info}
  each success in successes
    li(style="color: green") #{success}
  each error in errors
    li(style="color: red") #{error}
\end{lstlisting}
\risultato
\begin{browser}[localhost:3000/showflash]
  \begin{itemize}[leftmargin=1.8em, itemsep=0pt, topsep=2pt, parsep=0pt]
    \item {\color{css-blue}Sei stato reindirizzato.}
    \item {\color{css-blue}Questa è una dimostrazione dei messaggi flash.}
    \item {\color{css-green}È andato tutto bene.}
    \item {\color{css-red}Ma abbiamo anche un messaggio di errore.}
  \end{itemize}
\end{browser}
\wtnota{Ricaricando la pagina, la lista è vuota.}
\end{affianco}

\begin{console}[Terminale]
$ curl -i -c jar3 localhost:3000/redirect
HTTP/1.1 302 Found
X-Powered-By: Express
Location: /showflash
...
Set-Cookie: connect.sid=s%3Ap9hBc3nGrmeoUYJ6jF7H1CXS8J3nD17b.a0XjrBePGj5DidRQpfbEml0EPI7awuDpYX2SqTznca4; Path=/; HttpOnly
...

Found. Redirecting to /showflash
$ curl -b jar3 -c jar3 localhost:3000/showflash
<ul><li style="color: blue">Sei stato reindirizzato.</li><li style="color: blue">Questa è una dimostrazione dei messaggi flash.</li><li style="color: green">È andato tutto bene.</li><li style="color: red">Ma abbiamo anche un messaggio di errore.</li></ul>
$ curl -b jar3 -c jar3 localhost:3000/showflash
<ul></ul>
\end{console}

La prima risposta è un \code{302} con l'intestazione \code{Location}: il
browser la segue da solo, \code{curl} no (lo farebbe con l'opzione \code{-L}),
per questo la seconda richiesta è scritta a mano. Una volta consumati, i
messaggi flash vengono \textbf{cancellati}: ricaricando \code{/showflash}, o
visitandola direttamente senza essere stati reindirizzati, non compare alcun
messaggio. È precisamente il comportamento che si desidera per una notifica
\guillemotleft usa e getta\guillemotright.

\info{Un avviso di Node.js}{
  Con Node.js 22, all'avvio di un programma che usa \code{connect-flash} il
  terminale del server mostra \textit{DeprecationWarning: The `util.isArray`
  API is deprecated}. Il pacchetto è vecchio e usa una funzione di Node che
  verrà eliminata: per ora l'avviso è innocuo, e il pacchetto funziona.}

\section{Riferimenti}

\begin{itemize}
  \item \textbf{Express web framework (Node.js/JavaScript)} --- MDN web docs. \\
        \urlx{https://developer.mozilla.org/en-US/docs/Learn/Server-side/Express_Nodejs}
  \item \textbf{Inversion of Control}, Martin Fowler. \\
        \urlx{https://martinfowler.com/bliki/InversionOfControl.html}
  \item \textbf{Express Guide} --- sito ufficiale. \\
        \urlx{https://expressjs.com/en/guide/routing.html} \\
        \urlx{https://expressjs.com/en/guide/writing-middleware.html} \\
        \urlx{https://expressjs.com/en/guide/using-middleware.html} \\
        \urlx{https://expressjs.com/en/guide/using-template-engines.html} \\
        \urlx{https://expressjs.com/en/guide/error-handling.html}
  \item \textbf{Migrating to Express 5} --- sito ufficiale. \\
        \urlx{https://expressjs.com/en/guide/migrating-5.html}
\end{itemize}
